\documentclass[./main.tex]{subfiles}
\begin{document}

% Slide # 8
\begin{frame}[t, label=slide08]
        % Title
        \frametitle{Theorem n. 2 (linear stability)}

        % Body
        \footnotesize

        \begin{theorem}[Hartman-Grobman]\label{thm:hartman_grobman}
                Let $\dot{x}=f(x;\mu)$ be a smooth dynamical system and $x^{*}$ be a\alr{hyperbolic}equilibrium of $f$. Then there exist a neighborhood $U \subset \mathbb{R}^{n}$ of $x^{*}$ and a homeomorphism $h:U\to \mathbb{R}^{n}$ s.t. $h(\phi(x,t))=\psi(h(x),t)$, where $\psi$ is the solution stemming from the origin of the linearised system $\dot{x}=Df(x^{*})(x-x^{*})$.
        \end{theorem}

        The following vector field has three equilibria: $(0,0)$, $(1,0)$ and $(0,1)$
        
        \begin{equation*}
                \begin{cases}
                    \dot{x} = x(1-x-y)\,, \\
                    \dot{y} = y(2-x-2y)\,.
                \end{cases}
        \end{equation*}
        
        \begin{columns}[T]
                
                \column{0.33\textwidth}

                \vspace{-2mm}
                \begin{figure}[H]
                        \centering 
                        \includegraphics<2->[keepaspectratio, width=\textwidth]{./figures/slide_08_1.png}%
                \end{figure}

                \hfill

                \column{0.33\textwidth}

                \vspace{-2mm}
                \begin{figure}[H]
                        \centering 
                        \includegraphics<3->[keepaspectratio, width=\textwidth]{./figures/slide_08_2.png}%
                \end{figure}

                \hfill

                \column{0.33\textwidth}

                \vspace{-2mm}
                \begin{figure}[H]
                        \centering 
                        \includegraphics<4->[keepaspectratio, width=\textwidth]{./figures/slide_08_3.png}%
                \end{figure}

       \end{columns}

\end{frame}

\end{document}
